\documentclass[11pt]{article}
\usepackage{amsmath,amssymb,amsthm}
\newtheorem{proposition}{Proposition}
\title{An Exact Audit of Isolated Divided-Difference Tokens}
\author{}
\date{November 24, 2026}
\begin{document}
\maketitle

\paragraph{Scope.}
This is a written obstruction to a specific proposed decoder, not a
lower bound on attention head complexity. It rules out increasing the
polynomial degree in an isolated-token divided-difference construction
as a source of unbounded independent bit readout. Interacting token
assignments and permutation-dependent rank drops are not covered.

\begin{proposition}
Let $P$ be any real polynomial, let $D_0,D_1,\ldots,D_k$ be fixed real
numbers, and restrict the variable token $t$ to a set on which
$D_i-t>0$ for every $i$. Define
\[
 F(D,t)=\frac{P(D)-P(t)}{D-t},\qquad
 R_i(t)=F(D_i,t)-F(D_0,t).
\]
If the $k$ readouts $R_i$ jointly attain every strict sign word, then
$k\leq2$. More generally, if fixed real thresholds $b_i$ are allowed
and the readouts are $R_i(t)+b_i$, then $k\leq3$.
The conclusions also hold when the token is restricted to any finite
library, regardless of its size or the degree of $P$.
\end{proposition}

\begin{proof}
If $D_i=D_0$, then the unthresholded readout is identically zero,
and the thresholded readout is constant. Such a coordinate already
prevents all strict words. We may assume $D_i\ne D_0$.
Let $L_i(t)$ be the affine function whose graph is the secant of $P$
through $D_0$ and $D_i$:
\[
 L_i(t)=\frac{P(D_i)(t-D_0)+P(D_0)(D_i-t)}{D_i-D_0}.
\]
Direct common-denominator subtraction gives the exact identity
\[
 R_i(t)=
 \frac{(D_i-D_0)(P(t)-L_i(t))}{(D_i-t)(D_0-t)}.
 \tag{1}
\]
Both denominator factors are positive. Removing their product and
the fixed nonzero factor $D_i-D_0$ only applies a fixed sign flip to
coordinate $i$. The transformed unthresholded output vector therefore
lies in a fixed affine subspace of $\mathbb R^k$ of dimension at most
two, with scalar coordinates $t$ and $P(t)$.

For completeness, a proper affine subspace $a+V$ of $\mathbb R^k$
cannot meet every open orthant. Choose a nonzero vector $w$ orthogonal
to $V$. Every vector $z\in a+V$ satisfies $w\cdot z=w\cdot a$.
If this constant is nonnegative, the orthant with signs opposite to
the nonzero coordinates of $w$ makes $w\cdot z<0$. If the constant
is negative, choose the matching signs, which give $w\cdot z>0$.
Coordinates where $w$ vanishes may have arbitrary strict signs.
Thus all strict words require the affine subspace to have dimension
at least $k$, proving the first assertion.

For thresholded outputs, multiply by the same positive denominator
product and divide by $D_i-D_0$. The resulting readout is
\[
 P(t)-L_i(t)+\frac{b_i}{D_i-D_0}(D_i-t)(D_0-t).
 \tag{2}
\]
Its vector lies in a fixed affine subspace parametrized by
$t,t^2,P(t)$, of dimension at most three. The preceding orthant
argument proves the second assertion. Restricting the token set
cannot add sign words.
\end{proof}

\paragraph{Use in the present research.}
Summing common-probe divided differences over a permutation of a
fixed token library gives a permutation-invariant background. If a
query changes exactly one probe and subtracts that background, its
token-dependent contribution is precisely (1). Arbitrarily large
values or high-degree polynomial interpolation in the numerator do
not change this bound. A proposed multi-token construction must give
an explicit additional interaction; it cannot use this isolated
readout as an unbounded-bit primitive.

\paragraph{Verification boundary.}
The identity and affine-subspace proof above are written only. This
file adds no Lean declaration or model definition and claims no
asymptotic upper bound for $W(n)$.
\end{document}
